\section{Conclusion}\label{sec:conclusion}

Does hiring from cartels change competitive behavior? Yes---but in the opposite direction from what the contagion hypothesis predicts. Firms that absorb ex-cartel workers win more procurement auctions ($+5$--$6$~pp) without raising their prices. The effect is cartel-specific ($3.7\times$ larger than absorbing workers from non-cartel exits) and driven by operational workers, not by the managers and professionals who might carry coordination know-how. The mechanism is capacity transfer: ex-cartel workers bring product-market expertise---familiarity with specific items, supply chains, and procurement procedures---that helps the receiving firm compete more effectively in the markets the cartel once dominated.

\medskip

\noindent\textbf{Policy implications.} For competition authorities concerned about the unintended consequences of enforcement, the result is reassuring. Cartel dissolution triggers labor reallocation from convicted firms to competitors, but this reallocation strengthens competition rather than propagating collusion. Debarment policies need not worry that displaced workers will seed the next cartel---the workers who move are operational staff, and their effect on receiving firms is pro-competitive.

A separate lesson concerns econometric practice. The naive TWFE estimate produces a striking price ``contagion'' effect ($+0.26$ log points, $t > 5$) that is entirely an artifact of staggered-treatment bias. Had we not implemented the stacked estimator, the paper would have reported a false positive with direct policy implications---a cautionary tale for applied researchers working with staggered treatments.

\noindent\textbf{Limitations.} The principal caveat is that RAIS records employment bonds, not tasks. We observe that operational workers move from cartel firms to competitors, but we cannot observe what knowledge they carry or how the receiving firm deploys them. The win-rate increase could reflect the workers' direct contribution (executing contracts that the cartel firm can no longer fulfill) or an indirect effect (the receiving firm uses the new workers to bid more aggressively, knowing it can deliver). Distinguishing these channels requires richer data on within-firm task assignment.

The pre-trend at $\tau = -2$ ($t = 1.83$) in the stacked DiD specification, while below the 5\% threshold, is a caveat. The CEM-matched specification eliminates it entirely (max $|t| = 0.65$), suggesting selection on pre-period levels rather than a genuine anticipation effect. Future work with more cartels and more cohorts would provide sharper identification.

\noindent\textbf{Looking ahead.} The finding that enforcement redistributes operational capacity from cartels to competitors opens a new set of questions. Does the capacity transfer improve procurement efficiency---lower prices, better quality, faster delivery? Do the former cartel firms recover, or does the loss of workers mark a permanent competitive decline? And does the pattern generalize beyond procurement to product markets where cartels are enforced? These questions connect the enforcement literature to the broader economics of creative destruction, where firm exit and worker reallocation are the mechanisms through which markets restructure after competitive shocks.
